(a) A family of closed subsets with no minimal member yields a strictly descending chain. Conversely, a minimal member of a descending chain forces the chain to stabilize. Thus (i) and (ii) are equivalent. Taking complements gives (i)$\iff$(iii) and (ii)$\iff$(iv).

(b) The finite unions of members of an open cover have a maximal member by (a). It contains every member of the cover, so it is $X$ and gives a finite subcover.

(c) Write a descending chain of closed subsets of $Y\subseteq X$ as $Y\cap F_n$, with $F_n$ closed in $X$. Replacing $F_n$ by $F_1\cap\cdots\cap F_n$ makes the $F_n$ descend, so the chain stabilizes.

(d) A noetherian space has finitely many irreducible components. An irreducible Hausdorff space has at most one point, since two distinct points have disjoint nonempty open neighborhoods. Thus $X$ is finite, and a finite Hausdorff space is discrete.
